\section{Introduction}

Gauss's \textit{Theorema Egregium} (1827) states that the Gaussian curvature of a surface is an intrinsic invariant: it can be determined entirely from measurements within the surface (the first fundamental form), without reference to the ambient space. This was a revolutionary discovery that founded differential geometry.

In this chapter, we cover:
\begin{itemize}
    \item First and second fundamental forms
    \item Gaussian curvature and mean curvature
    \item The Theorema Egregium
    \item Geodesics and the Gauss-Bonnet theorem
    \item Numerical computation of curvature
\end{itemize}

\section{Surfaces and Fundamental Forms}

Let $\mathbf{r}(u,v)$ be a parametrized surface in $\R^3$.

\begin{definition}[First Fundamental Form (Metric Tensor)]
\[
I = E du^2 + 2F du dv + G dv^2,
\]
where
\[
E = \mathbf{r}_u \cdot \mathbf{r}_u,\quad
F = \mathbf{r}_u \cdot \mathbf{r}_v,\quad
G = \mathbf{r}_v \cdot \mathbf{r}_v.
\]
\end{definition}

\begin{definition}[Second Fundamental Form]
Let $\mathbf{N} = \frac{\mathbf{r}_u \times \mathbf{r}_v}{\|\mathbf{r}_u \times \mathbf{r}_v\|}$ be the unit normal.
\[
II = L du^2 + 2M du dv + N dv^2,
\]
where
\[
L = \mathbf{r}_{uu} \cdot \mathbf{N},\quad
M = \mathbf{r}_{uv} \cdot \mathbf{N},\quad
N = \mathbf{r}_{vv} \cdot \mathbf{N}.
\]
\end{definition}

\section{Curvature}

\begin{definition}[Gaussian and Mean Curvature]
\[
K = \frac{LN - M^2}{EG - F^2}, \quad
H = \frac{EN - 2FM + GL}{2(EG - F^2)}.
\]
\end{definition}

\begin{theorem}[Theorema Egregium]
The Gaussian curvature $K$ depends only on the first fundamental form $E, F, G$ and their derivatives. It is an intrinsic invariant of the surface.
\end{theorem}

\begin{proof}[Proof Sketch]
The Christoffel symbols $\Gamma_{ij}^k$ are computed from $E, F, G$. The Riemann curvature tensor $R_{ijkl}$ is expressed in terms of $\Gamma$. The Gaussian curvature is $K = R_{1212} / \det(g)$. No dependence on $L, M, N$ remains.
\end{proof}

\section{Christoffel Symbols and Riemann Tensor}

\begin{align*}
\Gamma_{11}^1 &= \frac{GE_u - 2FF_u + FE_v}{2(EG-F^2)} \\
\Gamma_{12}^1 &= \frac{GE_v - FG_u}{2(EG-F^2)} \\
\Gamma_{22}^1 &= \frac{2GF_v - GG_u - FG_v}{2(EG-F^2)} \\
\Gamma_{11}^2 &= \frac{2EF_u - EE_v - FE_u}{2(EG-F^2)} \\
\Gamma_{12}^2 &= \frac{EG_u - FE_v}{2(EG-F^2)} \\
\Gamma_{22}^2 &= \frac{EG_v - 2FF_v + FG_u}{2(EG-F^2)}
\end{align*}

Then $R_{1212} = \frac{1}{2}(E_{vv} + G_{uu} - 2F_{uv}) + \text{terms in } \Gamma$.

\section{Python Implementation}

In \texttt{python/theorema.py}:

\begin{lstlisting}[style=python, caption={Key functions in python/theorema.py}]
first_fundamental_form(r, u, v)     # Returns E, F, G
second_fundamental_form(r, u, v)    # Returns L, M, N
christoffel_symbols(E, F, G)        # Returns Gamma^k_ij
gaussian_curvature(E, F, G, ...)    # K from first form only
mean_curvature(E, F, G, L, M, N)    # H
riemann_tensor(E, F, G)             # R_ijkl
theorema_egregium_check(r, u, v)    # Verify K intrinsic
geodesic_equations(r)               # Geodesic ODEs
gauss_bonnet_discrete(vertices, faces)  # Discrete Gauss-Bonnet
\end{lstlisting}

\section{Numerical Examples}

\begin{example}[Sphere of Radius $R$]
$\mathbf{r}(\theta,\phi) = (R\sin\theta\cos\phi, R\sin\theta\sin\phi, R\cos\theta)$
$E = R^2$, $F = 0$, $G = R^2\sin^2\theta$, $K = 1/R^2$, $H = 1/R$.
\begin{lstlisting}[style=python]
>>> from theorema import gaussian_curvature, first_fundamental_form
>>> import numpy as np
>>> r = lambda u,v: np.array([np.sin(u)*np.cos(v), np.sin(u)*np.sin(v), np.cos(u)])
>>> E, F, G = first_fundamental_form(r, 1.0, 0.5)
>>> K = gaussian_curvature(r, 1.0, 0.5)
>>> print(f"K = {K:.6f}")  # 1.0
\end{lstlisting}
\end{example}

\begin{example}[Cylinder]
$\mathbf{r}(\theta,z) = (R\cos\theta, R\sin\theta, z)$
$K = 0$ (intrinsically flat), $H = 1/(2R)$.
\end{example}

\begin{example}[Pseudosphere]
Surface of revolution of the tractrix: $K = -1$ (constant negative curvature).
\end{example}

\section{Exercises}

\begin{exercise}
Prove that the Gaussian curvature of a surface of revolution $z = f(r)$ is $K = \frac{f' f''}{r(1+f'^2)^2}$.
\end{exercise}

\begin{exercise}
Implement discrete differential geometry operators on triangle meshes to compute $K$ at vertices.
\end{exercise}

\begin{exercise}
Verify the Theorema Egregium numerically: compute $K$ from first form only and compare with $K$ from both forms for various surfaces.
\end{exercise}

\begin{exercise}
Apply the discrete Gauss-Bonnet theorem to a triangulated torus and verify $\sum K_i A_i = 0$.
\end{exercise}

\section{References}

\begin{itemize}
    \item \cite{gauss-disquisitiones-generales} -- Gauss, \textit{Disquisitiones Generales circa Superficies Curvas} (1827)
    \item \cite{do-carmo} -- do Carmo, \textit{Differential Geometry of Curves and Surfaces}
    \item \cite{spivak-dg} -- Spivak, \textit{A Comprehensive Introduction to Differential Geometry}
    \item \cite{meyer-discrete} -- Meyer et al., \textit{Discrete Differential-Geometry Operators for Triangulated 2-Manifolds}
\end{itemize}